Use \(\mathrm{lem:positive\text{-}multistate\text{-}native\text{-}forest\text{-}separator}\) and put
\[
 h(a)=P_{\overline M_0}(R=a\mid\operatorname{do}(X_N\cap T=0))
 -P_{\overline M_1}(R=a\mid\operatorname{do}(X_N\cap T=0)),
\]
with \(\delta=h(r^\dagger)\ne0\). For each \(r\in R\), native incidence supplies a path
\[
 \gamma_r:r=v_{r,0}\to\cdots\to v_{r,\ell_r}=y(r),
 \qquad y(r)\in Y_N,
\]
inside \(G^S[N_S\setminus X_N]\). Paths may share vertices or branches.

Augment each path vertex by one signal coordinate for every path through it. The \(r\)-signal has alphabet \(\mathcal X_r\). At its source and each later path vertex use the same randomized-response channel
\[
 K_{\epsilon,r}(z\mid u)
 =(1-\epsilon)\mathbf 1\{z=u\}+\frac{\epsilon}{|\mathcal X_r|},
 \qquad 0<\epsilon<1.
\]
Use independent finite private noises, and product channels at shared vertices. Retain the Kivva base coordinates on \(T\), and give all other base coordinates independent full-support laws. Conditional on the base vector on \(T\), all new coordinates have one common strictly positive conditional kernel. Equality and positivity of the base laws therefore imply
\[
 P_0(V=v)=P_1(V=v)>0
\]
for every augmented state \(v\).

Give all vertices in \(A_S\) identical independent full-support mechanisms, and make \(S\) an independent fair bit while ignoring its declared parents. Then
\[
 P_j(V=v,S=1)=\frac12P_j(V=v)>0,\qquad
 P_j(V=v\mid S=1)=P_j(V=v),
\]
so the selected laws agree.

Let \(L=\sum_{r\in R}(\ell_r+1)\). Couple every channel as exact copying with probability \(1-\epsilon\) and uniform redrawing with probability \(\epsilon\). All terminal signals equal their root inputs with probability at least \((1-\epsilon)^L\). If \(A_Y\) is the event that the terminal signal tuple equals \(r^\dagger\), then for \(j=0,1\),
\[
 \left|P_j(A_Y\mid\operatorname{do}(X=x^\dagger))
 -P_{\overline M_j}(R=r^\dagger\mid\operatorname{do}(X_N\cap T=0))\right|
 \le1-(1-\epsilon)^L.
\]
Here \(x^\dagger\) uses Kivva's zero levels on \(X_N\cap T\) and fixed values elsewhere. The paths avoid \(X_N\), so intervention cuts no signal channel. Choose \(\epsilon\) with
\[
 2\{1-(1-\epsilon)^L\}<|\delta|.
\]
The triangle inequality yields
\[
 |P_0(A_Y\mid\operatorname{do}(X=x^\dagger))
 -P_1(A_Y\mid\operatorname{do}(X=x^\dagger))|
 \ge|\delta|-2\{1-(1-\epsilon)^L\}>0.
\]
Independence of \(S\) persists after intervention, so conditioning on \(S=1\) changes neither target.

All used forest and routing edges are actual edges of \(G^S\); the premodels induce subgraphs of \(G^S\). Apply \(\operatorname{Exactify}^{\mathrm{fin}}_{G^S}\). By \(\mathrm{lem:finite\text{-}exactification\text{-}preserves\text{-}laws}\), the completed models induce exactly \(G^S\) and preserve the equal positive selected law and target contrast. They lie in \(\mathfrak M_+(G^S,\widetilde\sigma)\). This one finite-signature pair falsifies \(\mathrm{def:graphwise\text{-}sid}\).